% 10-modern-methods — the section HBP wrote as one paragraph and one figure.
% Subsection D ends the section completely: the equal-time cross-check, if it
% arrives, is added as its own file (10e) referenced from nowhere else.
\section{Comparison with modern methods}
\label{sec:modern}

The 1990 paper compared against the two external results then available:
't~Hooft's large-$N$ solution and Hamer's finite-lattice SU(2)
spectrum~\cite{Hamer:1982nc}.  The modern record is far richer, and this
section is organized by what each method can check.

\subsection{Large $N$ and integrability}
\label{sec:modern-largeN}

The 't~Hooft equation is now solved to essentially arbitrary precision: its
equivalence to a TQ-Baxter system~\cite{Ambrosino:2023dik,Litvinov:2024aa}
turns the large-$N$ meson spectrum into a benchmark exact for our purposes.
Our finite-$N$ masses, rescaled to 't~Hooft's normalization
(Appendix~\ref{app:units}), approach these values with the expected $1/N^2$
systematics, updating Fig.~8(b) of Ref.~\cite{Hornbostel:1988fb} with error
bars on both axes' inputs.  The same rescaled comparison exposes what
large $N$ misses at strong coupling --- light baryons, and the U(1)
distinction of Sec.~\ref{sec:massless} --- exactly as the original argued,
now with the residuals quantified.

\subsection{Lattice}
\label{sec:modern-lattice}

Hamer's SU(2) finite-lattice masses remain the only equal-time spectrum
across the full coupling range at $N = 2$.  We compare against his published
Table~1, whose stated uncertainties at weak coupling are of order the values
themselves; a digitization of his figure would suggest a precision his own
table disclaims, and we deliberately do not use one.  Within those bars,
agreement is complete, including the $N = 2$ identity $M_{\rm bar} = M_{\rm
mes}$ that his data was the first to exhibit.  Modern overlap-fermion
lattice work confirms the expected infrared behaviour at
$(N_c, N_f) = (2,1)$ and $(3,1)$~\cite{Karthik:2023vjv}, though without
spectrum tables to compare digit-for-digit.

\subsection{Tensor networks and quantum hardware}
\label{sec:modern-tn}

Hamiltonian-lattice calculations with matrix product states now reach
SU(2)~\cite{Silvi:2016cas,Banuls:2017ena} and
SU(3)~\cite{Silvi:2019wnf,Hayata:2023bgh} in one dimension, and variational
quantum hardware has produced SU(2) meson and baryon
masses~\cite{Atas:2021ext} with SU(3) preparations
following~\cite{Farrell:2022wyt}.  These are staggered-fermion,
finite-lattice-spacing computations, so the honest comparison is at matched
continuum extrapolations rather than raw numbers; where the published works
quote continuum values, we tabulate them beside ours.  Precision benchmarks
in the Schwinger model~\cite{Banuls:2013jaa,Dempsey:2022nys} calibrate what
tensor-network methods achieve at U(1), and our equal-time
exact-diagonalization cross-checks (Appendix~\ref{app:data}) reproduce them
in our own conventions.  No public code exists for any of the SU($N$)
tensor-network calculations above; the open pipeline released with this
paper is, to our knowledge, the first reproducible finite-$N$ spectrum code
in either quantization.

\subsection{The chiral-exponent tension}
\label{sec:modern-tension}

One discrepancy in this landscape is structural, and we state it as the open
question it is.  Every light-front determination of the finite-$N$ chiral
exponent --- this work, lightcone conformal truncation~\cite{Anand:2021qnd},
and $1/N$-corrected 't~Hooft perturbation theory~\cite{Kochergin:2024aa} ---
finds $M^2 \propto (m/g)^\alpha$ with $\alpha \to 1$, while the one
equal-time measurement, an MPS study of SU(2) over $m/g \in
[0.1, 0.4]$~\cite{Banuls:2017ena}, reports $\nu = 0.700(29)(11)$ for
$M \propto (m/g)^{\nu}$, i.e.\ $\alpha \simeq 1.4$: consistent with the
bosonization prediction and several standard deviations from the light-front
value.  The methods share no assumption more specific than quantization
itself --- but they differ in precisely the one all light-front schemes
share, the trivial vacuum, which is what an anomalous chiral exponent is
about.  The equal-time side of the ledger is not yet decisive either: the
measured window spans a factor of four in $m/g$, not a decade, and no
SU($N$) analogue of the additive staggered mass shift~\cite{Dempsey:2022nys}
--- the same order in the lattice spacing as the mass being fitted --- has
been published or applied.  Resolving the tension requires an equal-time
computation in the same chiral window with that renormalization under
control; it is beyond the published record, and beyond this paper.
